\documentclass[11pt]{article}

\usepackage{amsmath,amssymb,amsthm,mathtools}
\usepackage[margin=1.05in]{geometry}
\usepackage{hyperref}

\newtheorem{theorem}{Theorem}
\newtheorem{lemma}{Lemma}
\newtheorem{corollary}{Corollary}
\newtheorem{definition}{Definition}
\newtheorem{remark}{Remark}

\newcommand{\E}{\mathbb{E}}
\newcommand{\N}{\mathbb{N}}
\newcommand{\PhiH}{\Phi_H}

\title{Even Girth Gives Infinitely Many Tensor--Tur\'{a}n Violations}
\author{}
\date{}

\begin{document}

\maketitle

\section{Statement and conventions}

All finite graphs in this note are simple.  Let \(H\) be a finite graph,
and put
\[
    n=v(H)=|V(H)|,
    \qquad
    m=e(H)=|E(H)|.
\]
Its chromatic polynomial is denoted by \(\chi_H\).  If \(H\) contains a
cycle, its \emph{girth} \(g\) is the minimum length of a cycle in \(H\),
and \(c_g(H)\) is the number of edge sets of \(H\) that form a cycle of
length \(g\).

A graphon is a symmetric measurable function
\[
    W\colon \Omega^2\longrightarrow[0,1]
\]
on a probability space \((\Omega,\mu)\).  Its edge density and its
homomorphism densities are
\[
    p(W):=t(K_2,W)=\int_{\Omega^2}W(x,y)\,d\mu(x)d\mu(y)
\]
and
\[
    t(F,W)
    :=
    \int_{\Omega^{V(F)}}
       \prod_{uv\in E(F)}W(x_u,x_v)
       \prod_{u\in V(F)}d\mu(x_u).
\]
This probability-space formulation is equivalent to the usual
\([0,1]\)-formulation for the graphons used below.  In particular, a
finite product of copies of \([0,1]\) is an atomless standard probability
space and can be transported to \([0,1]\) by a measure-space isomorphism,
without changing any homomorphism density.

For \(0\le p<1\), define the target
\[
    \PhiH(p)
    :=
    (1-p)^n
    \chi_H\!\left(\frac{1}{1-p}\right).
\]

The purpose of this note is to prove the following statement.  It is
phrased directly as a violation of the target inequality, rather than
only as non-optimality of a Tur\'{a}n graphon.

\begin{theorem}[Even-girth obstruction]\label{thm:main}
Let \(H\) be a finite simple graph of finite even girth.  There exists an
integer \(k_0=k_0(H)\ge \max\{2,\chi(H)-1\}\) such that, for every integer
\(k\ge k_0\), there are an integer \(Q\ge1\) and a graphon \(W_{k,Q}\)
such that
\[
    p(W_{k,Q})=1-\frac1k
    \ge
    1-\frac{1}{\chi(H)-1}
\]
and
\[
    t(H,W_{k,Q})
    <
    \PhiH\!\left(1-\frac1k\right).
\]
Consequently the proposed inequality fails for \(H\) at infinitely many
edge densities in its required density interval.
\end{theorem}

No numerical search is used.  The graphons \(W_{k,Q}\) are explicit
finite tensor products of balanced Tur\'{a}n graphons.

\section{Balanced Tur\'{a}n graphons}

\begin{definition}[Balanced Tur\'{a}n graphon]
For an integer \(q\ge2\), partition \([0,1]\) into measurable intervals
\(I_1,\dots,I_q\), each of measure \(1/q\), and set
\[
    T_q(x,y)
    :=
    \begin{cases}
        0,&x,y\in I_i\text{ for some }i,\\
        1,&x\in I_i,\ y\in I_j,\ i\ne j.
    \end{cases}
\]
\end{definition}

\begin{lemma}[Densities in \(T_q\)]\label{lem:turan-density}
For every finite simple graph \(F\) and every integer \(q\ge2\),
\[
    p(T_q)=1-\frac1q,
    \qquad
    t(F,T_q)=\frac{\chi_F(q)}{q^{v(F)}}.
\]
In particular,
\[
    t(H,T_q)=\PhiH\!\left(1-\frac1q\right).
\]
\end{lemma}

\begin{proof}
For independent uniform points \(X,Y\in[0,1]\), the two points lie in
different parts with probability \(1-1/q\).  This proves the first
identity.  For a family of independent uniform points
\((X_u)_{u\in V(F)}\), let \(a(u)\in[q]\) be the index of the part
containing \(X_u\).  The vector \(a\colon V(F)\to[q]\) is uniform among
all \(q^{v(F)}\) maps, and
\[
    \prod_{uv\in E(F)}T_q(X_u,X_v)=1
\]
exactly when \(a\) is a proper \(q\)-colouring of \(F\).  Thus the
integral defining \(t(F,T_q)\) is
\(\chi_F(q)/q^{v(F)}\).  The last identity follows by substituting
\(1-p=1/q\) into the definition of \(\Phi_H\).
\end{proof}

Define, for integers \(q\ge2\),
\[
    F_H(q)
    :=
    \frac{t(H,T_q)}{p(T_q)^m}
    =
    \frac{\chi_H(q)}{q^n}
    \left(1-\frac1q\right)^{-m}.
\]
Equivalently,
\[
    t(H,T_q)
    =
    \left(1-\frac1q\right)^m F_H(q).
\]

\section{The first non-forest coefficient}

The next proof is expanded down to finite-set and polynomial statements
that can be formalized independently of asymptotic graph theory.

For \(A\subseteq E(H)\), let \(c(A)\) be the number of connected
components of the spanning graph \((V(H),A)\), including isolated
vertices, and define
\[
    r(A):=n-c(A).
\]

\begin{lemma}[Whitney expansion]\label{lem:whitney}
For every positive integer \(q\),
\[
    \chi_H(q)
    =
    \sum_{A\subseteq E(H)}(-1)^{|A|}q^{n-r(A)}.
\]
Hence the equality holds as a polynomial identity in \(q\).
\end{lemma}

\begin{proof}
For each edge \(e=uv\), let \(B_e\) be the event that a map
\(f\colon V(H)\to[q]\) satisfies \(f(u)=f(v)\).  Proper colourings are
the maps outside \(\bigcup_{e\in E(H)}B_e\).  Inclusion--exclusion gives
\[
    \chi_H(q)
    =
    \sum_{A\subseteq E(H)}(-1)^{|A|}
      \left|\bigcap_{e\in A}B_e\right|.
\]
A map in the intersection is constant on each connected component of
\((V(H),A)\), and these component values are arbitrary.  The cardinality
of the intersection is therefore \(q^{c(A)}=q^{n-r(A)}\).  This proves
the formula for every positive integer \(q\); two integer polynomials
agree identically when they agree at infinitely many integers.
\end{proof}

\begin{lemma}[Rank classification below the girth]\label{lem:rank-classification}
Assume \(H\) has finite girth \(g\).  Then:
\begin{enumerate}
    \item if \(0\le j\le g-2\) and \(r(A)=j\), then \(A\) is a forest,
    so \(|A|=j\);
    \item if \(r(A)=g-1\), then either \(A\) is a forest with
    \(|A|=g-1\), or \(A\) is exactly the edge set of one \(g\)-cycle.
\end{enumerate}
\end{lemma}

\begin{proof}
Rank is monotone under inclusion.  A cycle of length \(\ell\) has rank
\(\ell-1\).  Thus, if \(A\) is cyclic and \(C\subseteq A\) is a cycle,
then
\[
    r(A)\ge r(C)=|C|-1\ge g-1.
\]
This proves the first assertion, since every forest has rank equal to its
number of edges.

Now suppose \(A\) is cyclic and \(r(A)=g-1\).  The displayed inequalities
force every cycle \(C\subseteq A\) to have length \(g\) and
\(r(C)=r(A)\).  Fix one such \(C\), and consider the connected components of
the spanning graph \((V(H),A)\).  From the definition \(r(A)=n-c(A)\), its
rank is the sum of \(|V(D)|-1\) over its connected components \(D\).  The
component containing the \(g\) vertices
of \(C\) already contributes \(g-1\) to the rank.  If it contained any
additional nonisolated vertex, then its contribution would be at least \(g\).
Any other component containing an edge would contribute at least \(1\).
Both alternatives contradict \(r(A)=g-1\).  Therefore every edge of \(A\)
has both endpoints in \(V(C)\).  An edge in \(A\setminus E(C)\) would then join
two nonconsecutive vertices of \(C\).  It is a chord and, together with either
of the two paths between its endpoints on \(C\), creates a cycle shorter than
\(g\).  This contradicts the definition of girth.  Hence \(A=E(C)\).
\end{proof}

\begin{lemma}[First correction term]\label{lem:first-correction}
If \(H\) has finite girth \(g\), then there is a constant \(C_H\ge0\)
such that, for every sufficiently large integer \(q\),
\[
    \left|
      F_H(q)-1-
      (-1)^g\frac{c_g(H)}{(q-1)^{g-1}}
    \right|
    \le
    \frac{C_H}{(q-1)^g}.
\]
In particular, if \(g\) is even, then \(F_H(q)>1\) for every sufficiently
large integer \(q\).
\end{lemma}

\begin{proof}
Group the Whitney expansion by rank.  By
Lemma~\ref{lem:rank-classification}, the coefficient of \(q^{n-j}\) is
\[
    (-1)^j\binom{m}{j}
    \qquad(0\le j\le g-2),
\]
because every rank-\(j\) set is a \(j\)-edge forest.  At rank \(g-1\),
the forests contribute
\[
    (-1)^{g-1}\binom{m}{g-1},
\]
and the \(c_g(H)\) girth cycles contribute
\[
    (-1)^g c_g(H).
\]
On the other hand, the binomial theorem gives
\[
    q^n\left(1-\frac1q\right)^m
    =
    \sum_{j=0}^{m}(-1)^j\binom{m}{j}q^{n-j}.
\]
Consequently there is a polynomial \(R_H(q)\) of degree at most \(n-g\)
such that
\[
    \chi_H(q)
    =
    q^n\left(1-\frac1q\right)^m
    +(-1)^g c_g(H)q^{n-g+1}
    +R_H(q).
\]
(If the displayed degree bound is negative, \(R_H=0\).)

Because \(R_H\) has finitely many coefficients, there is \(C_0\ge0\)
such that \(|R_H(q)|\le C_0q^{n-g}\) for \(q\ge1\).  Dividing by
\(q^n(1-1/q)^m\), and using
\((1-1/q)^{-m}\le2^m\) for \(q\ge2\), gives a constant \(C_1\) with
\[
    F_H(q)
    =
    1+(-1)^g c_g(H)q^{-(g-1)}+\varepsilon_q,
    \qquad
    |\varepsilon_q|\le C_1q^{-g}.
\]
The exact identity
\[
    q^{-(g-1)}
    =(q-1)^{-(g-1)}
      \left(1-\frac1q\right)^{g-1}
\]
and the finite binomial expansion of
\((1-1/q)^{g-1}\) show that
\[
    \left|q^{-(g-1)}-(q-1)^{-(g-1)}\right|
    \le C_2(q-1)^{-g}
\]
for a constant \(C_2\) depending only on \(g\).  Enlarging the constant
proves the asserted error bound.

If \(g\) is even, then \(c_g(H)\ge1\).  For every \(q\) with
\(q-1>C_H/c_g(H)\), the positive leading term is strictly larger than
the absolute value of the error.  Therefore \(F_H(q)>1\).
\end{proof}

\section{Tensor products and the telescoping construction}

\begin{definition}[Tensor product]
If \(U\colon\Omega_1^2\to[0,1]\) and
\(W\colon\Omega_2^2\to[0,1]\) are graphons, define their tensor product
on the product probability space by
\[
    (U \otimes W)((x_1,x_2),(y_1,y_2))
    :=U(x_1,y_1)W(x_2,y_2).
\]
\end{definition}

\begin{lemma}[Tensor multiplicativity]\label{lem:tensor}
For every finite graph \(F\),
\[
    p(U \otimes W)=p(U)p(W),
    \qquad
    t(F,U \otimes W)=t(F,U)t(F,W).
\]
The same identities hold for every finite iterated tensor product.
\end{lemma}

\begin{proof}
Insert the definition of \(U \otimes W\) into the relevant integral.
The integrand becomes the product of one function of all
\(\Omega_1\)-coordinates and one function of all
\(\Omega_2\)-coordinates.  Tonelli's theorem applies because the
integrand is nonnegative and bounded.  The integral over the product
space therefore factors into the two defining homomorphism-density
integrals.  The edge identity is the special case \(F=K_2\), and finite
iteration follows by induction on the number of factors.
\end{proof}

\begin{lemma}[A finite-product estimate]\label{lem:product-estimate}
Let \(a_{Q,1},\dots,a_{Q,Q}\) be real numbers.  If
\[
    S_Q:=\sum_{j=1}^{Q}|a_{Q,j}|\longrightarrow0,
\]
then
\[
    \prod_{j=1}^{Q}(1+a_{Q,j})\longrightarrow1.
\]
\end{lemma}

\begin{proof}
Expanding the product and taking absolute values term by term gives
\[
    \left|\prod_{j=1}^{Q}(1+a_{Q,j})-1\right|
    \le
    \prod_{j=1}^{Q}(1+|a_{Q,j}|)-1.
\]
Since \(1+x\le e^x\) for \(x\ge0\), the right-hand side is at most
\(e^{S_Q}-1\), which tends to zero.
\end{proof}

\begin{lemma}[Telescoping tensor products]\label{lem:telescoping}
Fix an integer \(k\ge2\).  For every integer \(Q\ge1\), set
\[
    W_{k,Q}
    :=
    \bigotimes_{j=1}^{Q} T_{Q(k-1)+j}.
\]
If \(H\) has finite girth, then
\[
    p(W_{k,Q})=1-\frac1k
\]
for every \(Q\), and
\[
    \lim_{Q\to\infty}t(H,W_{k,Q})
    =
    \left(1-\frac1k\right)^m.
\]
\end{lemma}

\begin{proof}
By Lemma~\ref{lem:tensor} and the definition of \(T_q\),
\[
\begin{aligned}
    p(W_{k,Q})
    &=
    \prod_{j=1}^{Q}
       \left(1-\frac{1}{Q(k-1)+j}\right)\\
    &=
    \prod_{a=Q(k-1)+1}^{Qk}\frac{a-1}{a}\\
    &=
    \frac{Q(k-1)}{Qk}
    =1-\frac1k.
\end{aligned}
\]
This is an exact telescoping identity for each \(Q\).

Again by tensor multiplicativity,
\[
\begin{aligned}
    t(H,W_{k,Q})
    &=
    \prod_{j=1}^{Q}t\!\left(H,T_{Q(k-1)+j}\right)\\
    &=
    \left(1-\frac1k\right)^m
      \prod_{j=1}^{Q}F_H(Q(k-1)+j).
\end{aligned}
\]
Lemma~\ref{lem:first-correction} implies that there are constants
\(C\) and \(M\), depending only on \(H\), such that
\[
    |F_H(a)-1|\le Ca^{-(g-1)}
    \qquad(a\ge M).
\]
For all sufficiently large \(Q\), every
\(a=Q(k-1)+j\) is at least \(M\), and
\[
\begin{aligned}
    \sum_{j=1}^{Q}|F_H(Q(k-1)+j)-1|
    &\le
    C\sum_{j=1}^{Q}(Q(k-1)+j)^{-(g-1)}\\
    &\le
    C Q\,[Q(k-1)]^{-(g-1)}\\
    &=
    C(k-1)^{-(g-1)}Q^{2-g}.
\end{aligned}
\]
Every finite simple graph has \(g\ge3\), so the last expression tends
to zero.  Lemma~\ref{lem:product-estimate} now shows that the product of
the \(F_H\)-terms tends to \(1\), proving the claimed limit.
\end{proof}

\section{Proof of the obstruction theorem}

\begin{proof}[Proof of Theorem~\ref{thm:main}]
Let \(g\) be the even girth of \(H\).  By
Lemma~\ref{lem:first-correction}, there is \(k_1\) such that
\[
    F_H(k)>1
    \qquad(k\ge k_1).
\]
Set
\[
    k_0:=\max\{2,k_1,\chi(H)-1\}.
\]
Fix \(k\ge k_0\).  Lemma~\ref{lem:turan-density} and the definition of
\(F_H\) give
\[
\begin{aligned}
    \PhiH\!\left(1-\frac1k\right)
    &=t(H,T_k)\\
    &=\left(1-\frac1k\right)^mF_H(k)\\
    &>\left(1-\frac1k\right)^m.
\end{aligned}
\]
By Lemma~\ref{lem:telescoping}, every \(W_{k,Q}\) has edge density
\(1-1/k\), and its \(H\)-density converges to the strictly smaller last
quantity.  By the definition of convergence, there is \(Q\ge1\) for
which
\[
    t(H,W_{k,Q})
    <
    \PhiH\!\left(1-\frac1k\right).
\]
Finally, \(k\ge\chi(H)-1>0\) implies
\[
    1-\frac1k
    \ge
    1-\frac{1}{\chi(H)-1}.
\]
Thus the violation lies in the required density interval.  Since this
works for every integer \(k\ge k_0\), it gives infinitely many distinct
edge densities.
\end{proof}

\begin{corollary}
Every non-bipartite finite graph of even girth violates the proposed
chromatic-polynomial lower bound at infinitely many admissible edge
densities.  The same construction also applies to bipartite graphs of
finite even girth; non-bipartiteness is not used in the proof.
\end{corollary}

\section{Formalization dependency map}

In the integrated Lean project, this still-unformalized public theorem will
belong in \path{lean/Taeyoung/Methods/EvenGirthTensor/Main.lean}.  With the
common predicates from \path{lean/Taeyoung/Foundation/Status.lean}, its
intended interface is:

\begin{verbatim}
theorem evenGirth_violatesLowerBound
    (hgirth : HasFiniteEvenGirth H)
    (heven : Even hgirth.girth) :
    ViolatesLowerBound H
\end{verbatim}

Here \texttt{ViolatesLowerBound H} packages an admissible edge density, an
explicit graphon certificate, and a strict comparison with the shared
polynomially continued chromatic target.  Thus an Atlas example imports the
method's \texttt{Main.lean}, verifies the finite graph hypothesis, and applies
one theorem; it does not reproduce the asymptotic or tensor calculation.
The project-level organization is specified in
\path{LEAN_VERIFICATION_ARCHITECTURE.md}.

The proof above deliberately separates the ingredients that a Lean
verification must supply.

\begin{enumerate}
    \item \textbf{Finite graph layer.}  Define \(c(A)\), rank
    \(r(A)=n-c(A)\), forests, cycles, and girth.  Prove rank monotonicity,
    rank of a forest, and rank of a cycle.  Lemma~\ref{lem:rank-classification}
    is then finite combinatorics.
    \item \textbf{Chromatic-polynomial layer.}  Formalize finite
    inclusion--exclusion to prove Lemma~\ref{lem:whitney}.  Coefficient
    extraction below girth is a finite sum calculation; no analytic
    theorem is involved.
    \item \textbf{Elementary real analysis.}  Bound a fixed polynomial
    by a constant times its leading power, prove the displayed conversion
    from \(q^{-g+1}\) to \((q-1)^{-g+1}\), and use the
    epsilon definition of convergence.  Lemma~\ref{lem:product-estimate}
    avoids taking logarithms or assuming the factors are positive.
    \item \textbf{Measure layer.}  Homomorphism density is a finite
    product integral of bounded measurable functions.  Tensor
    multiplicativity is one application of Tonelli/Fubini followed by
    induction over the finite edge set.
    \item \textbf{Transport to \([0,1]\).}  Either define graphons over
    arbitrary probability spaces, as above, or transport each finite
    product space to \([0,1]\).  Only invariance of finite product
    integrals under a measure-preserving isomorphism is needed.
\end{enumerate}

There is no appeal to regularity, injective subgraph counts, minimum
degree, pointwise degree bounds, or a limiting infinite tensor product.
For each counterexample, \(Q\) is finite.

%%%%%%%%%%%%%%%%%%%%%%%%%%%%%%%%%%%%%%%%%%%%%%%%%%%%%%%%%%%%%%%%%%%%%%%%%%%%%%%
\appendix
\section{What the Lean formalization actually proves}
%%%%%%%%%%%%%%%%%%%%%%%%%%%%%%%%%%%%%%%%%%%%%%%%%%%%%%%%%%%%%%%%%%%%%%%%%%%%%%%

Nineteen catalogue rows are \texttt{negative} on the strength of this note, each
with a \texttt{ViolatesLowerBound} theorem in
\texttt{lean/Taeyoung/Methods/Negative/}, over
\texttt{Methods/Negative/Tensor.lean}.  One difference of substance.

\subsection*{Theorem~\ref{thm:main} is not formalized; nineteen witnesses are}

The theorem asserts the \emph{existence} of $k_0$ and $Q$, and its proof is
asymptotic: Lemmas~\ref{lem:whitney}, \ref{lem:rank-classification},
\ref{lem:first-correction}, \ref{lem:product-estimate} and
\ref{lem:telescoping} together produce a $k$ beyond which the first correction
term dominates.  None of that is formalized.  What the catalogue needs is not
existence but a witness, and every row here is settled by a single explicit
\emph{two}-factor product $T_a\otimes T_b$ --- $Q=2$ throughout --- found by
that asymptotic analysis and then checked directly.

For such a witness the whole argument is the identity
\[
 t(H,T_a\otimes T_b)=\frac{\chi_H(a)\,\chi_H(b)}{(ab)^{v(H)}},
 \qquad
 p=\Bigl(1-\frac1a\Bigr)\Bigl(1-\frac1b\Bigr),
\]
which is Lemma~\ref{lem:turan-density} together with Lemma~\ref{lem:tensor},
followed by exact rational arithmetic against $\PhiH(p)$.  Both $\chi_H(a)$ and
$\chi_H(b)$ are evaluated as counts of proper assignments by kernel reduction,
so no chromatic polynomial has to be known in advance.

The tensor product itself is built without any measure-theoretic product.  On
$\mathrm{Fin}\,a\times\mathrm{Fin}\,b$ with the uniform measure two points are
adjacent exactly when they differ in both coordinates, homomorphism density is
a finite sum, and an assignment is proper for the product exactly when both of
its coordinate projections are proper --- so the factorization is
$\mathtt{Equiv.arrowProdEquivProdArrow}$ and a product of sums.  The transport
to $[0,1]$ discussed at the end of Section~5 is therefore not needed either.

\end{document}
